\begin{claim}{D.51}{Proved}{The second index-three kernel and its three peripheral normalizers}
Let $a$ have order three in the full punctured-torus
loop-outer group, put $F=\widehat F(x,y)$,
$c=[x,y]$, and let $E$ be its centre-free extension.
Besides $F$, there is exactly one other torsion-free
normal subgroup $U$ of index three in $E$.
Write $V=F\cap U$. Then
\[
 \begin{gathered}
 [F:V]=[U:V]=3,\qquad E/V\cong C_3^2,\\
 V=\widehat{\pi_1(S_{1,3})},\qquad
 U\ \hbox{is projective},\qquad
 U^{(3)}\cong\widehat F^{(3)}_2 .
 \end{gathered}                                         \tag{D.51a}
\]
The cusp normalizer $N=N_E(\langle c\rangle)$
is contained in $U$. The three cusp subgroups
$L_0,L_1,L_2$ of the displayed degree-three
cover have normalizers $N_i=N_U(L_i)$ with
the exact sequences
\[
              1\longrightarrow L_i
                 \longrightarrow N_i
                 \longrightarrow U/V\longrightarrow1.
                                                        \tag{D.51b}
\]
They are three distinct $U$-conjugacy classes,
cyclically permuted by $E/U$. All have the
given multiplier $\lambda$ under the same
generator label in $U/V$.

For a prime $q\ne3$ at which $\lambda_q$ is
a nonidentity cube root of unity, the
$U/V$-character multiplicities of
$H_1(V,\Q_q)$, ordered as $(1,\lambda_q,\lambda_q^2)$,
are exactly one of
\[
                         (1,3,0),\qquad (0,2,2).
                                                        \tag{D.51c}
\]
Correspondingly $U^{(q)}$ is free pro-$q$ of rank
one or zero. These are constraints on the actual
subgroup $U$; they do not identify it with a full
free profinite group.
\end{claim}
\begin{proof}
\emph{The section and cusp lines in the actual abelian quotient.}
Use D.27 to choose an actual order-three representative
$\alpha$ and write $E=F\rtimes_\alpha\langle s\rangle$.
Let $u:E\to\F_3$ be its quotient character, with $u(s)=1$.
D.48 gives exactly two full section classes,
$S_0=\langle s\rangle$ and $S_1$, and gives
$E^{\rm ab,(3)}\cong C_3^2$.

Let $M$ be the reduction of the homology matrix of
$\alpha$ modulo three. By D.48 it is a nonidentity
unipotent matrix with $(M-I)^2=0$.
On $\F_3^2$ write
$\omega((a,b),(a',b'))=ab'-ba'$.
Choose $w$ with $e=(M-I)w\ne0$.
Then $(e,w)$ is a basis, $Me=e$, $Mw=w+e$,
and $d=\omega(e,w)\ne0$.
The linear functional
\[
                  \chi(r)=\omega(e,r)/d
                                                        \tag{D.51d}
\]
is $M$-invariant. It gives a character of $F$,
which we extend to $E$ by $\chi(s)=0$.
Thus $(u,\chi):E\to\F_3^2$ is the actual
abelian quotient above.

The generator of $S_1$ labelled by $u=1$
has $\chi$-value $b_1\ne0$.
Here is the required separation, using the maps
already proved in D.48 rather than an assumption
about section labels.
The Bockstein of $\chi$ has nonzero high periodic
multiples. On a section normalizer its restriction
is the Bockstein of $\chi$ restricted to the section:
the fixed-group first cohomology vanishes by D.48.
If $\chi$ vanished on both section classes,
all those multiples would have zero full Brown
restriction. This contradicts D.48f.
It follows that $b_1\ne0$.

To locate the cusp line, use the characteristic
class-two exponent-three quotient
\[
 Q=\F_3^2\times\F_3,\qquad
 (r,z)(r',z')=(r+r',z+z'+\tfrac12\omega(r,r')) .
                                                        \tag{D.51e}
\]
The images of $x,y$ are the two standard vectors
with central coordinate zero, and $c$ maps to
$Z=(0,1)$.
Since $\lambda_3=1$, the induced automorphism is
\[
               \theta(r,z)=(Mr,z+\ell(r))
                                                        \tag{D.51f}
\]
for a linear functional $\ell$.
For an element $(r,z)s$ of $Q\rtimes_\theta C_3$
with quotient label one, direct multiplication gives
\[
       ((r,z)s)^3
          =Z^{\,d\chi(r)^2+\ell(e)\chi(r)} .
                                                        \tag{D.51g}
\]
To check the formula, write $r=be+tw$, so
$\chi(r)=t$. The successive vector parts are
$r,r+te,r+2te$, whose sum is zero.
The sum of the three central correction terms
is $\ell((2I+M)r)=t\ell(e)$.
The three alternating cross terms, multiplied
by $1/2$, sum to $dt^2$ in $\F_3$.
This proves (D.51g) for every element, including
the arbitrary central coordinate $z$.

The image of the actual generator of $S_1$
has cube one. Equation (D.51g) and $b_1\ne0$
therefore give $\ell(e)=-db_1$.
Thus its cube polynomial is
\[
                           db(b-b_1).
                                                        \tag{D.51h}
\]
Choose $v\in N$ with $u(v)=1$, as in D.48.
It satisfies $v^3=c^\kappa$ with
$\kappa_3$ a unit. Its image in (D.51e)
has nonzero cube exponent. The two roots of
(D.51h) are $0$ and $b_1$; the only other
element of $\F_3$ is $-b_1$. Consequently
\[
             (u,\chi)(S_0)=\langle(1,0)\rangle,\quad
             (u,\chi)(S_1)=\langle(1,b_1)\rangle,\quad
             (u,\chi)(N)=\langle(1,-b_1)\rangle .
                                                        \tag{D.51i}
\]
The last equality uses that $N$ is generated
topologically by $v$ and $\langle c\rangle$,
and both characters kill $c$.

\emph{The full kernels and peripheral labels.}
Define the continuous characters and subgroups
\[
           \psi=\chi+b_1u,\qquad
           U=\ker\psi,\qquad V=\ker u\cap\ker\chi .
                                                        \tag{D.51j}
\]
These are actual subgroups of $E$.
The character $\psi$ is nonzero on both
section classes in (D.51i), and it kills $N$.
Every nonidentity finite subgroup of $E$
is one of the order-three section groups,
up to conjugacy. Hence $U$ is torsion-free.
The four lines in $E^{\rm ab,(3)}=\F_3^2$
give all normal index-three subgroups of $E$.
Two contain respectively $S_0$ and $S_1$.
The other two are exactly $\ker u=F$ and
$\ker\psi=U$, proving the uniqueness assertion.
Also $E=U\rtimes S_0$, since $\psi|_{S_0}$
is an isomorphism onto $\F_3$.

The restriction of $\chi$ to $F$ is nonzero
and kills its commutator cusp.
Choose a word $f$ in the dense discrete
group $F(x,y)$ with $\chi(f)=1$.
The degree-three covering of the punctured
torus corresponding to $\ker(\chi|_{F(x,y)})$
has three boundary lifts of degree one.
Its Euler characteristic is $-3$; thus its
type is $(1,3)$, and its fundamental group
has rank four.
The induced profinite subgroup is exactly $V$,
with cusp representatives
\[
          L_i=f^i\langle c\rangle f^{-i},
          \qquad
          N_i=f^i N f^{-i}\quad(0\leq i<3).
                                                        \tag{D.51k}
\]
The full peripheral normalizer assertion A.7
gives $N_E(L_i)=N_i$.
Since $N\subset U$ and $U$ is normal in $E$,
all $N_i$ lie in $U$. Their kernels under
$U\to U/V$ are exactly $L_i$, because
$N_i\cap F=L_i$; their images are all of
$U/V$, because $u(v)=1$.
This proves (D.51b), with the common generator
label supplied by $u$.

For the class assertions, the actual quotient
$(u,\chi)$ identifies $E/V$ with $\F_3^2$,
and the image of every $N_i$ is the same
line $U/V$.
The three $V$-conjugacy classes of the cusp
lifts are therefore the three cosets of this
line in $E/V$. The subgroup $U/V$ fixes each
class, so they remain distinct under $U$.
The quotient $E/U$ acts regularly on these
three cosets, proving its asserted cyclic
permutation. This also proves that conjugation
by a generator of $U/V$ fixes each cusp
class with multiplier $\lambda$.

The open subgroup $V$ of $U$ is free profinite
of rank four and $U$ is torsion-free.
Serre's finite-index cohomological-dimension
theorem, used in D.27, makes $U$ projective.
Its maximal pro-three quotient is projective
as a pro-three group, by the finite embedding
problem characterization, and is therefore
free pro-three.
The pro-three topology on $V$ is exactly the
one induced from $U$: taking the intersection
of the three conjugates of a normal
three-power-index subgroup of $V$ produces
one normal in $U$, still of three-power index.
Thus $V^{(3)}$ embeds as an open subgroup
of index three in $U^{(3)}$.
Schreier's formula now reads
\[
                 4=1+3(\operatorname{rank}U^{(3)}-1),
\]
and gives rank two. This proves (D.51a).

\emph{The exact character restriction at another prime.}
The three oriented boundary classes of $V$
have their sole relation given by their sum.
They form a two-dimensional submodule of
$H_1(V,\Q_q)$ on which $U/V$ acts by
$\lambda_q$.
Filling the three boundaries gives the compact
degree-three cover of the original torus.
Its map on homology is the inclusion of a
rank-two sublattice of index three, hence an
isomorphism over $\Q_q$, and over $\Z_q$
when $q\ne3$.
Conjugation by $v$ projects to the original
outer action $a$ on the base; its inner
factor acts trivially on that base homology.
Therefore, if $A_q$ denotes the original
matrix, the exact boundary sequence and
averaging over $C_3$ give
\[
       H_1(V,\Q_q)\cong
                  \lambda_q^{\oplus2}
                    \oplus(\Q_q^2,A_q)
                 \quad\hbox{as }U/V\hbox{-modules}.
                                                        \tag{D.51l}
\]
This comparison is induced by the stated
boundary inclusion and covering homology map.

At the specified prime, $A_q^3=I$ and
$\det A_q=\lambda_q$.
The three cube roots $1,\lambda_q,\lambda_q^2$
are distinct in $\Q_q$. Diagonalization gives
exactly the two eigenvalue multisets
\[
                 \{1,\lambda_q\},\qquad
                 \{\lambda_q^2,\lambda_q^2\}.
                                                        \tag{D.51m}
\]
Adding the two boundary characters in (D.51l)
proves (D.51c).
For $q\ne3$, the same decomposition holds
over $\Z_q$: the averaging idempotents have
unit denominator three, and the cube-root
differences are units.
Hochschild--Serre for $V\triangleleft U$ then
identifies $H^1(U,\F_q)$ with the invariant
subspace of $H^1(V,\F_q)$, of dimension one
or zero respectively.
The projectivity of $U$ makes its maximal
pro-$q$ quotient free pro-$q$, with this
number of generators. This proves the last
assertion.

The program
\path{code/part_c/torus_order3_sibling_d51.py}
checks (D.51g) for every central-linear
automorphism above every nonidentity unipotent
matrix over $\F_3$, the resulting section and
cusp lines, and the finite eigenvalue patterns
in (D.51m). These finite calculations validate
the displayed formulas. The full subgroup,
topology, and projectivity conclusions are
proved by the maps above.
\end{proof}

\paragraph{Exact remaining recognition problem.}
The construction supplies a specified torsion-free
projective group $U$ with an open full free subgroup
of rank four and three specified normalizers.
It supplies no identification of $U$ with
$\widehat F_2$.
At the primes in (D.51c), its maximal pro-$q$
quotient has rank zero or one, whereas the
maximal pro-three quotient has rank two.
The existence or exclusion of such a full
$U$ with these three normalizers remains the
step needed for this approach to determine
every torus order-three multiplier.
